% Claim proof file for row claim_3_16 -- "MarginalizationPreserves".
% Auto-generated stub by scaffold/solve_chapter.py at row start. A manager
% agent dedicated to writing the TeX proof fills in the proof body below.
%
% This file is a sibling of `main.tex` inside `<subsection>/tex/`. The
% `subfiles` package lets it render both standalone and as part of
% `main.tex` via `\subfile{...}`.
%
% The statement is restated above the proof so the file is self-contained
% when rendered on its own. The proof must:
%   - Stay close to the lecture notes' proof if one exists (always search
%     the LN first for a `\begin{proof}` block immediately following the
%     claim; copy / lightly edit when present).
%   - Otherwise use the LN paradigm and cite prior claims by their `ref`.
%   - Be self-contained enough that a mathematician can follow it without
%     re-reading the LN.
%   - Have no `sorry`, `omitted`, or "obvious" shortcuts.

\documentclass[main]{subfiles}

\begin{document}
% ref: claim_3_16 (proof, with statement restated for readability)
% title: MarginalizationPreserves
\def\rowref{claim\_3\_16}
\phantomsection\label{claim_3_16}

% --- statement (restated) ---------------------------------------------------
\begin{Rem}[MarginalizationPreserves]\label{rem:marg_preserves_ancestors_bifurcations_acyclicity}
Marginalization preserves ancestral relations, bifurcations and acyclicity:
    \begin{enumerate}
        \item For $v_1,v_2 \in G$ with $v_1,v_2 \notin W$ we have the equivalence:
            \[ v_1 \in \Anc^G(v_2)\quad \iff \quad v_1 \in \Anc^{G^{\sm W}}(v_2).   \]
          \item\label{rem:marg_preserves_bifurcations} 
            For $v_1,v_2 \in G \sm W$ (and, optionally, $v_3 \in G\sm W$): there is a bifurcation between $v_1$ and $v_2$ (with source $v_3$) in $G$ if and only if there is a bifurcation between $v_1$ and $v_2$ (with source $v_3$) in $G^{\sm W}$.
        \item If the CDMG $G$ is acyclic then so is $G^{\sm W}$ and a topological order of $G$ induces a
            topological order on $G^{\sm W}$ (by just ignoring the nodes from $W$).
    \end{enumerate}
\end{Rem}

% --- proof ------------------------------------------------------------------
\begin{proof}
% TODO: write the proof body.
\end{proof}

\end{document}
